
    \documentclass{scrreprt} 
    \usepackage{listings}
    \usepackage{underscore}
    \usepackage{graphicx}
    \usepackage[bookmarks=true]{hyperref}
    \usepackage[utf8]{inputenc}
    \usepackage[english]{babel}
    \hypersetup{
        bookmarks=false,    % show bookmarks bar?
        pdftitle={Software Requirement Specification},    % title
        pdfauthor={Jean-Philippe Eisenbarth},                     % author
        pdfsubject={TeX and LaTeX},                        % subject of the document
        pdfkeywords={TeX, LaTeX, graphics, images}, % list of keywords
        colorlinks=true,       % false: boxed links; true: colored links
        linkcolor=blue,       % color of internal links
        citecolor=black,       % color of links to bibliography
        filecolor=black,        % color of file links
        urlcolor=purple,        % color of external links
        linktoc=page            % only page is linked
    }% 
    \def\myversion{My Report}
    \date{}
    
    \begin{document}
    
    \begin{flushright}
        \rule{16cm}{5pt}\vskip1cm
        \begin{bfseries}
            \Huge{SOFTWARE REQUIREMENTS\ SPECIFICATION}\
            \vspace{1.5cm} 
            \\for\\
            \vspace{1.5cm}
            IICT WEBSITE\\
            \vspace{1.5cm}
            \LARGE{Version \myversion}\\
            \vspace{1.5cm}
            Prepared by : Muhammad Adnan\\ 
            \vspace{1.5cm}
            \today\
        \end{bfseries}
    \end{flushright}
    
    
    \tableofcontents
    
    \chapter{Introduction} 
    \section{Purpose}
    <p>This template illustrates organizing the functional requirements for the product by system features, the major services provided by the product. You may prefer to organize this section by use case, mode of operation, user class, object class, functional hierarchy, or combinations of these, whatever makes the most logical sense for your product.&gt;</p><h2>1.1&nbsp;&nbsp;&nbsp;System Feature 1</h2><p>&lt;Don’t really say “System Feature 1.” State the feature name in just a few words.&gt;</p><p>4.1.1&nbsp;&nbsp;&nbsp;&nbsp;Description and Priority</p><p>&lt;Provide a short description of the feature and indicate whether it is of High, Medium, or Low priority. You could also include specific priority component ratings, such as benefit, penalty, cost, and risk (each rated on a relative scale from a low of 1 to a high of 9).&gt;</p><p>4.1.2&nbsp;&nbsp;&nbsp;&nbsp;Stimulus/Response Sequences</p><p>&lt;List the sequences of user actions and system responses that stimulate the behavior defined for this feature. These will correspond to the dialog elements associated with use cases.&gt;</p><p>4.1.3&nbsp;&nbsp;&nbsp;&nbsp;Functional Requirements</p><p>&lt;Itemize the detailed functional requirements associated with this feature. These are the software capabilities that must be present in order for the user to carry out the services provided by the feature, or to execute the use case. Include how the product should respond to anticipated error conditions or invalid inputs. Requirements should be concise, complete, unambiguous, verifiable, and necessary. Use “TBD” as a placeholder to indicate when necessary information is not yet available.&gt;</p><p><br></p><p>&lt;Each requirement should be uniquely identified with a sequence number or a meaningful tag of some kind.&gt;</p><p>&nbsp;</p><p>REQ-1:&nbsp;&nbsp;&nbsp;</p><p>REQ-2:&nbsp;&nbsp;&nbsp;</p><h2>1.2&nbsp;&nbsp;&nbsp;System Feature 2 (and so on)</h2>
    \section{Document Conventions} 
    This template illustrates organizing the functional requirements for the product by system features, the major services provided by the product. You may prefer to organize this section by use case, mode of operation, user class, object class, functional hierarchy, or combinations of these, whatever makes the most logical sense for your product.>

\subsection*{1.1   System Feature 1}

<Don’t really say “System Feature 1.” State the feature name in just a few words.>

4.1.1    Description and Priority

<Provide a short description of the feature and indicate whether it is of High, Medium, or Low priority. You could also include specific priority component ratings, such as benefit, penalty, cost, and risk (each rated on a relative scale from a low of 1 to a high of 9).>

4.1.2    Stimulus/Response Sequences

<List the sequences of user actions and system responses that stimulate the behavior defined for this feature. These will correspond to the dialog elements associated with use cases.>

4.1.3    Functional Requirements

<Itemize the detailed functional requirements associated with this feature. These are the software capabilities that must be present in order for the user to carry out the services provided by the feature, or to execute the use case. Include how the product should respond to anticipated error conditions or invalid inputs. Requirements should be concise, complete, unambiguous, verifiable, and necessary. Use “TBD” as a placeholder to indicate when necessary information is not yet available.>

<Each requirement should be uniquely identified with a sequence number or a meaningful tag of some kind.>

REQ-1:

REQ-2:

\subsection*{1.2   System Feature 2 (and so on)}
    \section{Intended Audience and Reading Suggestions}
    This template illustrates organizing the functional requirements for the product by system features, the major services provided by the product. You may prefer to organize this section by use case, mode of operation, user class, object class, functional hierarchy, or combinations of these, whatever makes the most logical sense for your product.>

\subsection*{1.1   System Feature 1}

<Don’t really say “System Feature 1.” State the feature name in just a few words.>

4.1.1    Description and Priority

<Provide a short description of the feature and indicate whether it is of High, Medium, or Low priority. You could also include specific priority component ratings, such as benefit, penalty, cost, and risk (each rated on a relative scale from a low of 1 to a high of 9).>

4.1.2    Stimulus/Response Sequences

<List the sequences of user actions and system responses that stimulate the behavior defined for this feature. These will correspond to the dialog elements associated with use cases.>

4.1.3    Functional Requirements

<Itemize the detailed functional requirements associated with this feature. These are the software capabilities that must be present in order for the user to carry out the services provided by the feature, or to execute the use case. Include how the product should respond to anticipated error conditions or invalid inputs. Requirements should be concise, complete, unambiguous, verifiable, and necessary. Use “TBD” as a placeholder to indicate when necessary information is not yet available.>

<Each requirement should be uniquely identified with a sequence number or a meaningful tag of some kind.>

REQ-1:

REQ-2:

\subsection*{1.2   System Feature 2 (and so on)}
    \section{Product Scope}
    This template illustrates organizing the functional requirements for the product by system features, the major services provided by the product. You may prefer to organize this section by use case, mode of operation, user class, object class, functional hierarchy, or combinations of these, whatever makes the most logical sense for your product.>

\subsection*{1.1   System Feature 1}

<Don’t really say “System Feature 1.” State the feature name in just a few words.>

4.1.1    Description and Priority

<Provide a short description of the feature and indicate whether it is of High, Medium, or Low priority. You could also include specific priority component ratings, such as benefit, penalty, cost, and risk (each rated on a relative scale from a low of 1 to a high of 9).>

4.1.2    Stimulus/Response Sequences

<List the sequences of user actions and system responses that stimulate the behavior defined for this feature. These will correspond to the dialog elements associated with use cases.>

4.1.3    Functional Requirements

<Itemize the detailed functional requirements associated with this feature. These are the software capabilities that must be present in order for the user to carry out the services provided by the feature, or to execute the use case. Include how the product should respond to anticipated error conditions or invalid inputs. Requirements should be concise, complete, unambiguous, verifiable, and necessary. Use “TBD” as a placeholder to indicate when necessary information is not yet available.>

<Each requirement should be uniquely identified with a sequence number or a meaningful tag of some kind.>

REQ-1:

REQ-2:

\subsection*{1.2   System Feature 2 (and so on)}
    \section{References} 
    This template illustrates organizing the functional requirements for the product by system features, the major services provided by the product. You may prefer to organize this section by use case, mode of operation, user class, object class, functional hierarchy, or combinations of these, whatever makes the most logical sense for your product.>

\subsection*{1.1   System Feature 1}

<Don’t really say “System Feature 1.” State the feature name in just a few words.>

4.1.1    Description and Priority

<Provide a short description of the feature and indicate whether it is of High, Medium, or Low priority. You could also include specific priority component ratings, such as benefit, penalty, cost, and risk (each rated on a relative scale from a low of 1 to a high of 9).>

4.1.2    Stimulus/Response Sequences

<List the sequences of user actions and system responses that stimulate the behavior defined for this feature. These will correspond to the dialog elements associated with use cases.>

4.1.3    Functional Requirements

<Itemize the detailed functional requirements associated with this feature. These are the software capabilities that must be present in order for the user to carry out the services provided by the feature, or to execute the use case. Include how the product should respond to anticipated error conditions or invalid inputs. Requirements should be concise, complete, unambiguous, verifiable, and necessary. Use “TBD” as a placeholder to indicate when necessary information is not yet available.>

<Each requirement should be uniquely identified with a sequence number or a meaningful tag of some kind.>

REQ-1:

REQ-2:

\subsection*{1.2   System Feature 2 (and so on)}
    
    \chapter{Overall Description} 
    \section{Product Perspective}
    <h2>1.1&nbsp;&nbsp;&nbsp;System Feature 1</h2><p>&lt;Don’t really say “System Feature 1.” State the feature name in just a few words.&gt;</p><p>4.1.1&nbsp;&nbsp;&nbsp;&nbsp;Description and Priority</p><p>&lt;Provide a short description of the feature and indicate whether it is of High, Medium, or Low priority. You could also include specific priority componednt ratings, such as benefit, penalty, cost, and risk (each rated on a relative scale from a low of 1 to a high of 9).&gt;</p>
    \section{Product Functions}
    <h2>1.1&nbsp;&nbsp;&nbsp;System Feature 1</h2><p>&lt;Don’t really say “System Feature 1.” State the feature name in just a few words.&gt;</p><p>4.1.1&nbsp;&nbsp;&nbsp;&nbsp;Description and Priority</p><p>&lt;Provide a short description of the feature and indicate whether it is of High, Medium, or Low priority. You could also include specific priority component ratings, such as benefit, penalty, cost, and risk (each rated on a relative scale from a low of 1 to a high of 9).&gt;d</p>
    \section{User Classes and Characteristics}
    <h2>1.1&nbsp;&nbsp;&nbsp;System Feature 1</h2><p>&lt;Don’t really say “System Feature 1.” State the feature name in just a few words.&gt;</p><p>4.1.1&nbsp;&nbsp;&nbsp;&nbsp;Description and Priority</p><p>&lt;Provide a short description of the feature and indicate whether it is of High, Medium, or Low priority. You could also include specific priority component ratings, such as benefit, penalty, cost, and risk (each rated on a relative scale from a low of 1 to a high of 9).&gt;</p>
    \section{Operating Environment}
    <h2>1.1&nbsp;&nbsp;&nbsp;System Feature 1</h2><p>&lt;Don’t really say “System Feature 1.” State the feature name in just a few words.&gt;</p><p>4.1.1&nbsp;&nbsp;&nbsp;&nbsp;Description and Priority</p><p>&lt;Provide a short description of the feature and indicate whether it is of High, Medium, or Low priority. You could also include specific priority component ratings, such as benefit, penalty, cost, and risk (each rated on a relative scale from a low of 1 to a high of 9).&gt;</p>
    \section{Design and Implementation Constraints}
    <h2>1.1&nbsp;&nbsp;&nbsp;System Feature 1</h2><p>&lt;Don’t really say “System Feature 1.” State the feature name in just a few words.&gt;</p><p>4.1.1&nbsp;&nbsp;&nbsp;&nbsp;Description and Priority</p><p>&lt;Provide a short description of the feature and indicate whether it is of High, Medium, or Low priority. You could also include specific priority component ratings, such as benefit, penalty, cost, and risk (each rated on a relative scale from a low of 1 to a high of 9).&gt;</p>
    \section{User Documentation}
    <h2>1.1&nbsp;&nbsp;&nbsp;System Feature 1</h2><p>&lt;Don’t really say “System Feature 1.” State the feature name in just a few words.&gt;</p><p>4.1.1&nbsp;&nbsp;&nbsp;&nbsp;Description and Priority</p><p>&lt;Provide a short description of the feature and indicate whether it is of High, Medium, or Low priority. You could also include specific priority component ratings, such as benefit, penalty, cost, and risk (each rated on a relative scale from a low of 1 to a high of 9).&gt;</p>
    \section{Assumptions and Dependencies}
    <h2>1.1&nbsp;&nbsp;&nbsp;System Feature 1</h2><p>&lt;Don’t really say “System Feature 1.” State the feature name in just a few words.&gt;</p><p>4.1.1&nbsp;&nbsp;&nbsp;&nbsp;Description and Priority</p><p>&lt;Provide a short description of the feature and indicate whether it is of High, Medium, or Low priority. You could also include specific priority component ratings, such as benefit, penalty, cost, and risk (each rated on a relative scale from a low of 1 to a high of 9).&gt;</p>
     
    \chapter{Features} 
    This template illustrates organizing the functional requirements for the product by system features, the major services provided by the product. You may prefer to organize this section by use case, mode of operation, user class, object class, functional hierarchy, or combinations of these, whatever makes the most logical sense for your product.>

\subsection*{1.1   System Feature 1}

<Don’t really say “System Feature 1.” State the feature name in just a few words.>

4.1.1    Description and Priority

<Provide a short description of the feature and indicate whether it is of High, Medium, or Low priority. You could also include specific priority component ratings, such as benefit, penalty, cost, and risk (each rated on a relative scale from a low of 1 to a high of 9).>

4.1.2    Stimulus/Response Sequences

<List the sequences of user actions and system responses that stimulate the behavior defined for this feature. These will correspond to the dialog elements associated with use cases.>

4.1.3    Functional Requirements

<Itemize the detailed functional requirements associated with this feature. These are the software capabilities that must be present in order for the user to carry out the services provided by the feature, or to execute the use case. Include how the product should respond to anticipated error conditions or invalid inputs. Requirements should be concise, complete, unambiguous, verifiable, and necessary. Use “TBD” as a placeholder to indicate when necessary information is not yet available.>

<Each requirement should be uniquely identified with a sequence number or a meaningful tag of some kind.>

REQ-1:

REQ-2:

\subsection*{1.2   System Feature 2 (and so on)}
      
    \chapter{Other Nonfunctional Requirements}
    \section{Performance Requirements} 
    If there are performance requirements for the product under various circumstances, state them here and explain their rationale, to help the developers understand the intent and make suitable design choices. Specify the timing relationships for real time systems. Make such requirements as specific as possible. You may need to state performance requirements for individual functional requirements or features. 
    \section{Saftey Requirements}
    Specify those requirements that are concerned with possible loss, damage, or harm that could result from the use of the product. Define any safeguards or actions that must be taken, as well as actions that must be prevented. Refer to any external policies or regulations that state safety issues that affect the product’s design or use. Define any safety certifications that must be satisfied. 
    \section{Security Requirements}
    Specify any requirements regarding security or privacy issues surrounding use of the product or protection of the data used or created by the product. Define any user identity authentication requirements. Refer to any external policies or regulations containing security issues that affect the product. Define any security or privacy certifications that must be satisfied. 
    \section{Software Quality Attributes}
    Specify any additional quality characteristics for the product that will be important to either the customers or the developers. Some to consider are: adaptability, availability, correctness, flexibility, interoperability, maintainability, portability, reliability, reusability, robustness, testability, and usability. Write these to be specific, quantitative, and verifiable when possible. At the least, clarify the relative preferences for various attributes, such as ease of use over ease of learning.>  
    \section{Business Rules} 
    List any operating principles about the product, such as which individuals or roles can perform which functions under specific circumstances. These are not functional requirements in themselves, but they may imply certain functional requirements to enforce the rules
    \end{document}